% ============================================================
\section{Discussão, limitações e conclusão}\label{sec:discussao}
% ============================================================

\paragraph{Três regiões do modelo.} Juntas, as medidas descrevem o Qwen3-4B-Base em três
regiões. Nos primeiros blocos (até o 5, aproximadamente), o vetor de uma ocorrência ainda é o
seu token: a rede do bloco 1 é a rede lexical com os empates desfeitos, e a linha mais próxima
em $E$ é a do próprio token. Nos blocos do meio (de 6 a 22), o vetor deixa de se parecer com o
próprio token e ainda não se parece com o próximo; é aí que as comunidades seguem o artigo, o
parágrafo e o tema, e que os sentidos de uma palavra ambígua mais se afastam. No último terço,
o vetor se aproxima do token que vem depois, as comunidades passam a segui-lo, a anisotropia
cresce e surgem \emph{hubs} fortes. A passagem entre as duas últimas regiões é abrupta: nos
blocos 22 a 25, a troca de vizinhos entre blocos, a concordância com o contexto e a distância
ao próximo token mudam juntas. A divisão é coerente com a ``lente de
\emph{logits}''~\cite{nostalgebraist2020}, em que as previsões só se formam nas camadas finais.

\paragraph{O que as redes mostram além das médias.} A autossimilaridade, a medida clássica de
contextualidade~\cite{ethayarajh2019}, cai até o meio do modelo e volta a subir no fim, igual
para palavras ambíguas e estáveis; sozinha, ela sugeriria que o último bloco é ``menos
contextual'' que o do meio. As redes mostram que a troca de vizinhos continua até o último
bloco, que as comunidades mudam de natureza (do texto para o próximo token) e que a separação
das palavras ambíguas cresce até o fim; a subida da autossimilaridade acompanha a anisotropia.
Mostram também que a identidade do token é persistente: a linha do token em $E$ deixa de ser a
mais próxima já no bloco 5, mas as ocorrências continuam vizinhas muito depois, numa direção que
não é mais o \emph{embedding} de entrada. E a fragmentação do português pesa: continuações de
palavra perdem a identidade logo, inícios de palavra a mantêm.

\paragraph{Limitações.}
\begin{itemize}
    \item \textbf{Os temas são recortes} próximos da categoria raiz, com alguns artigos pouco
    típicos (prêmios em Física, canções em Futebol). No núcleo, 39\% dos vértices vêm de
    categorias extras.
    \item \textbf{Os sentidos foram anotados por um modelo de linguagem}, sem revisão humana.
    \item \textbf{O Qwen3 fragmenta o português}: muitos vértices são pedaços de palavra, e as
    palavras-alvo ficaram restritas às que são um token só (\emph{corrente}, \emph{onda},
    \emph{célula} e \emph{núcleo} saíram por isso).
    \item \textbf{A amostra é desenhada}: as cotas (50 e 20 ocorrências) explicam parte da
    estrutura da rede lexical, e o núcleo
    domina as métricas globais.
    \item \textbf{Um modelo, uma língua, um tipo de texto}, que o modelo provavelmente viu no
    treinamento; os vetores carregam o arredondamento do \emph{bfloat16}. Não sabemos se as
    três regiões aparecem nos mesmos blocos em outros modelos.
\end{itemize}

\paragraph{Conclusão.} Redes k-NN sobre as mesmas ocorrências, uma por camada, permitem
acompanhar a reorganização das representações de um Transformer vértice a vértice e em várias
escalas, e o objetivo geral e as quatro perguntas foram respondidos. No Qwen3-4B-Base e em
texto enciclopédico em português, a reorganização é grande, não monótona e dependente do tipo
de token: a identidade lexical cede espaço aos poucos sem desaparecer, o contexto do texto
domina o meio do modelo, o próximo token domina o fim, e as palavras ambíguas separam as
ocorrências pelo sentido. Do lado da entrada, a rede do vocabulário é organizada por significado
entre línguas, e não por escrita.
